%% kapa: exam=2 round=20 year=2009 period=1 subject=appraisal_practice subject_ko=감정평가실무 seq=3 number=3 points=20
%% source: Q-Net 감정평가사 기출문제 (cst00309) artlSeq=5024180 file=extracted/1008/5024180/2009년 제20회 감정평가사 2차시험 기출문제/1교시.hwp
%% conversion: kordoc 4.24.51 -> tools/gen2.py (auto); review=proofread-sonnet; flags=none
\begin{longquestion}{【문제 3】}
감정평가사 S씨는 투자자로부터 부실채권(Non Performing Loan) 투자와 관련한 자문을 요청받았다. 부실채권은 당해부동산과 관련된 담보부채권이다. 주어진 자료를 활용하여 다음 물음에 답하시오. (20점) \par
\medskip\noindent (1) 가격시점 현재 대상부동산의 가격을 평가하시오. (10점) \par
\medskip\noindent (2) 대상부동산의 예상낙찰가를 낙찰가율과 낙찰사례를 통하여 각각 구해 결정하고, 법원의 경매절차 진행시 낙찰을 통해 대상부실채권으로부터 얻을 수 있는 예상현금흐름을 구하시오. (단, 시간적 요인은 고려하지 아니함) (10점) \par
<자료 1> 기본적 사항 \par
1) 대상부동산 \par
① 토지: A시 B구 C동 77번지, 대, 250m\textsuperscript{2}, 주거용, 일반상업지역, 세로에 접함, 장방형, 평지 \par
② 건물: 위 지상 벽돌조 슬래브지붕 2층건, 연면적 200m\textsuperscript{2} (1, 2층 각 100m\textsuperscript{2}) \par
2) 가격시점: 2009.09.06 \par
3) 개요 \par
① 대상부동산이 속한 A시 B구는 구도심내 일반상업지역인 C동, 아파트가 많이 소재하는 D동, 정비된 주택지대인 E동, 기타 F동 등으로 형성되어 있으며, 대상부동산의 주변은 구도심내 일반상업지역으로 로변으로는 다소 노후화된 3～4층 규모의 상업용건물이 소재하고 후면으로는 노후화된 주상용건물, 주거용건물 등이 혼재하여 있음. 도심지 재개발과 관련하여 사업을 추진중인 추진위원회는 설립되어 있으나 구체적인 계획은 미정인 상태임. \par
② 본건은 노후화된 2층의 주거용건물로써 1층에는 소유자가 거주하고 있으며 2층 일부는 임차인이 거주하고 있음. \par
③ 본건 주변의 거래상황은 재개발가능성을 염두에 둔 수요가 다소 있어 매도호가는 다소 상승중인 것으로 조사되었으며, 거래관행은 본건주변건물이 대체로 노후화되어 있어 토지면적만을 기준으로 가격이 형성되어 있는 것으로 조사됨. \par
4) 대상부실채권(NPL) \par
상기 대상부동산에 관련된 M은행의 500,000,000원의 담보부채권으로서 2순위로 근저당 설정되어 있음. (미납이자 등은 고려하지 아니함) \par
<자료 2> 표준지 공시지가(공시기준일 2009년 1월 1일) \par
본건과 가장 비교가능성 있는 다음의 표준지를 기준함. \par
{\small\centering
\begin{tabular}{|c|c|c|c|c|c|c|c|c|}
\hline
\makecell{일련번호} & \makecell{소재지} & \makecell{면적(m\textsuperscript{2})} & \makecell{지목} & \makecell{이용상황} & \makecell{용도지역} & \makecell{도로\newline{}교통} & \makecell{형상 및 지세} & \makecell{공시지가(원/m\textsuperscript{2})} \\
\hline
\makecell{1} & \makecell{A시 B구\newline{}C동 78} & \makecell{260} & \makecell{대} & \makecell{주거용} & \makecell{일반상업} & \makecell{세로} & \makecell{정방형\newline{}평지} & \makecell{5,000,000} \\
\hline
\end{tabular}\par}

<자료 3> 거래사례 \par
1) 토지: C동 100번지, 대, 300m\textsuperscript{2}, 주거용, 일반상업지역, 소로에 접함, 사다리형, 완경사 \par
2) 건물: 벽돌조 슬래브지붕 2층건, 연면적 200m\textsuperscript{2} \par
3) 거래가격: 1,455,000,000원 \par
4) 거래시점: 2009.05.01 \par
5) 기타: 본거래에 특이사항은 없었던 것으로 판단됨. \par
<자료 4> 임대내역 등 \par
1) 본건 1층은 소유자 자가사용이고, 2층 일부는 임차인에게 임대중이나 정확한 내역은 미상임. \par
2) 주변 탐문조사 결과 본건을 임대할 경우 1, 2층 각각 보증금 50,000,000원, 월세 1,400,000원에 임대가능한 것으로 조사됨. \par
3) 제경비는 임차인 부담으로 필요제경비, 공익비 및 실비초과액 등은 고려하지 아니함. \par
<자료 5> 낙찰사례 \par
1) 토지: C동 60번지, 대, 350m\textsuperscript{2}, 주거용, 일반상업지역, 세로에 접함, 사다리형, 완경사 \par
2) 건물: 벽돌조 슬래브지붕 2층, 연면적 180m\textsuperscript{2} \par
3) 경매평가금액(최초법사가): 1,600,000,000원 \par
4) 낙찰가: 1,070,000,000원 \par
5) 낙찰시점: 2009.06.01 \par
6) 기타: 경매당시 소유자와 일부 임차인이 거주 중이었고, 권리관계 등 제반사항은 본건과 유사한 것으로 조사되었음. \par
<자료 6> 토지개별요인 비교 \par
1) 도로: 세로(95), 소로(100), 중로(105), 광로(115) \par
2) 형상: 정방형(100), 장방형(100), 기타(95) \par
3) 지세: 평지(100), 완경사(95) \par
4) 기타: 토지의 기타 개별요인은 대상부동산과 표준지․사례들이 유사함. \par
<자료 7> 건물에 관한 사항 \par
1) 건물개요 \par
{\small\centering
\begin{tabular}{|c|c|c|c|}
\hline
\makecell{구분} & \makecell{대상건물} & \makecell{거래사례} & \makecell{낙찰사례} \\
\hline
\makecell{준공일자} & \makecell{1980.01.01} & \makecell{1982.01.01} & \makecell{1981.01.01} \\
\hline
\makecell{대지면적(m\textsuperscript{2})} & \makecell{250} & \makecell{300} & \makecell{350} \\
\hline
\makecell{연면적(m\textsuperscript{2})} & \makecell{200} & \makecell{200} & \makecell{180} \\
\hline
\makecell{구조} & \makecell{벽돌조\newline{}슬래브지붕} & \makecell{벽돌조\newline{}슬래브지붕} & \makecell{벽돌조\newline{}슬래브지붕} \\
\hline
\end{tabular}\par}

2) 벽돌조 슬래브지붕 건물신축단가(2009.01.01 기준): 700,000(원/m\textsuperscript{2}) \par
3) 내용연수 50년, 잔존가치 0\% \par
4) 감가상각은 만년감가함. \par
<자료 8> 낙찰가율 자료: 최근 6개월간 A시 B구 낙찰가율 \par
{\small\centering
\begin{tabular}{|c|c|}
\hline
\makecell{구 분} & \makecell{낙찰가율(\%)} \\
\hline
\makecell{아파트} & \makecell{80} \\
\hline
\makecell{단독주택} & \makecell{70} \\
\hline
\makecell{연립, 다세대 주택} & \makecell{68} \\
\hline
\makecell{상업용건물} & \makecell{73} \\
\hline
\makecell{기타} & \makecell{65} \\
\hline
\end{tabular}\par}

<자료 9> 시점수정 자료 \par
1) 지가변동률 \par
① 2009.01.01 ～ 가격시점: 1.00300 \par
② 2009.05.01 ～ 가격시점: 1.01000 \par
③ 2009.06.01 ～ 가격시점: 1.00700 \par
2) 건축비지수 \par
2009년 1월 1일 이후 건축비는 보합세임. \par
<자료 10> \par
현금흐름 산정시 검토할 이해관계는 다음과 같음. \par
1) 등기부상 \par
① 1순위 근저당(I은행): 400,000,000원 \par
② 2순위 근저당(M은행: 대상부실채권): 500,000,000원 \par
③ 3순위 근저당(N은행): 100,000,000원 \par
④ 4순위 근저당(P은행): 50,000,000원 \par
2) 기타 \par
① 경매감정평가 수수료 및 경매집행비용: 7,000,000원 \par
② 소액임차인: 16,000,000원 \par
③ 일반채권: 10,000,000원 \par
<자료 11> 기타사항 \par
1) 같은 동에서 소재하는 부동산은 동일한 지역요인을 가지는 것으로 조사됨. \par
2) 보증금 운용이율과 적용환원이율은 6\%로 함. \par
3) 근저당과 관련한 미납이자나 채권최고액 등은 고려하지 아니함. \par
\end{longquestion}
